\section{Lattice gauge theory and the transfer matrix}\label{sec:lattice}

Write $e_0,\dots,e_3$ for the unit vectors of $\Z^4$, $(x,\mu)$ for the link
from $x$ to $x+e_\mu$, and $Q_{(x,\mu)}^{-1}$ for the variable of the reversed
link. A box $\Lambda$ is a product of four finite intervals, with its set
$E(\Lambda)$ of links and its set $P(\Lambda)$ of plaquettes. For a plaquette
at $x$ in the plane $\mu<\nu$ with links $e_1=(x,\mu)$, $e_2=(x+e_\mu,\nu)$,
$e_3=(x+e_\nu,\mu)$, $e_4=(x,\nu)$, put $Q_p=Q_{e_1}Q_{e_2}Q_{e_3}^*Q_{e_4}^*$.
With $dQ$ the Haar probability measure of $G$, the Wilson measure of a box
with free boundary is
\[
W_\Lambda(dQ)=\frac1{Z_\Lambda}\exp\Bigl(N\beta\sum_{p\in P(\Lambda)}\ReTr
Q_p\Bigr)\prod_{e\in E(\Lambda)}dQ_e ,
\]
and the Wilson measure $\mu_L$ of the torus $\Lambda_L=(\Z/L\Z)^4$ is defined
in the same way with all links and plaquettes of $\Lambda_L$. Both are
invariant under the gauge transformations $Q_{(x,y)}\mapsto g_xQ_{(x,y)}
g_y^{-1}$. A function of finitely many links is a local observable, and it
is gauge invariant if it is invariant under these transformations. For
$\beta\ge0$ the measures are reflection positive for reflections in lattice
hyperplanes through sites and in hyperplanes bisecting links, on boxes
symmetric under the reflection and on tori of even side~\cite{OSei78}.

Let $\omega$ be a state on the gauge-invariant local observables of $\Z^4$
that is invariant under translations and reflection positive in this sense.
The Osterwalder--Schrader construction~\cite{OS73,OS75,OSei78} gives a
Hilbert space $\mathcal H$ with a unit vector $\Omega$, a map $F\mapsto\Psi_F$
from the observables at nonnegative times to $\mathcal H$ with dense range,
and a self-adjoint transfer matrix $0\le\mathbb T\le1$ with $\mathbb T\Omega=
\Omega$ and
\[
\ip{\Psi_F}{\mathbb T^t\Psi_H}=\omega\bigl(\vartheta(\bar F)\,\tau_tH\bigr),
\qquad t\in\Z_{\ge0},
\]
where $\vartheta$ is the time reflection and $\tau_t$ the time translation;
$\Psi_F$ is the class of $F$ for the positive semidefinite form
$\omega(\vartheta(\bar F)H)$, and $\|\Psi_F\|^2\le\omega(|F|^2)$. The
\emph{spectral gap} of $\mathbb T$ is $m$ if its spectrum on $\Omega^\perp$
lies in $[0,e^{-m}]$.

\begin{proposition}\label{prop:gapdecay}
Let $m>0$ and let $\mathcal F$ be a class of gauge-invariant local
observables, supported at nonnegative times, that is dense in $L^2(\omega)$
on each bounded time support. If for every $F\in\mathcal F$ there is
$C_F<\infty$ with
\[
\omega(\vartheta(\bar F)\tau_tF)-|\omega(F)|^2\le C_Fe^{-mt}\qquad(t\ge0),
\]
then $\mathbb T$ has spectral gap at least $m$.
\end{proposition}

\begin{proof}
Put $\Psi'_F=\Psi_F-\ip\Omega{\Psi_F}\Omega$, so that the left side equals
$\ip{\Psi'_F}{\mathbb T^t\Psi'_F}=\int\lambda^t\,d\sigma_F(\lambda)$, with
$\sigma_F$ the spectral measure of $\mathbb T$ at $\Psi'_F$. For $\delta>0$
this is at least $\sigma_F((e^{-m+\delta},1])e^{(-m+\delta)t}$, which exceeds
$C_Fe^{-mt}$ for large $t$ unless $\sigma_F((e^{-m+\delta},1])=0$. So the
spectral projection of $(e^{-m},1]$ annihilates every $\Psi'_F$. These vectors
are dense in $\Omega^\perp$, because $\|\Psi_F\|^2\le\omega(|F|^2)$ and the
class is dense, so the projection vanishes on $\Omega^\perp$.
\end{proof}

The decay of a single two-point function does not suffice: it places only
one spectral measure in $[0,e^{-m}]$, and the density of the class is what
gives the gap.

